\chapter{Literature Review}
\label{chap:lr}
\chaptermark{Second Chapter Heading}

This chapter provides a structured review of the existing research related to software architecture, maintainability, quality attributes, refactoring practices, technical debt, and architecture erosion.  
The objective of this review is to identify conceptual foundations, analytical methods, and empirical findings that support the investigation of how backend architectural styles influence software maintainability.

\section{Software Architecture Foundations}

Early research in software architecture established a conceptual level of reasoning that goes beyond individual components or algorithms. Software architecture is commonly defined as the organization of a system in terms of components, connectors, and their relationships, providing a high-level view of system structure~\cite{garlan1993introduction}. This perspective shifts design from isolated implementation decisions toward reusable architectural structures governed by explicit constraints.

Architectural styles such as Pipes-and-Filters, Layered Systems, Event-based architectures, and Repository-based architectures represent recurring structural solutions. Each style introduces specific invariants that restrict component interaction and shape system behavior. These structural constraints directly influence quality attributes such as modifiability, analyzability, and evolvability. In practice, many systems combine multiple architectural styles across different abstraction levels, which expands the architectural design space and introduces trade-offs between competing quality concerns~\cite{garlan1993introduction}.

In addition to structural organization, software architecture also captures design intent and rationale. A widely accepted model describes architecture as a combination of elements, their form, and the rationale behind design decisions~\cite{perry1992foundations}. This model emphasizes that architectural decisions are intentional and must be understood through multiple complementary views, including processing, data, and connector perspectives. Such multi-view descriptions are essential for understanding how architectural decisions affect system qualities over time~\cite{perry1992foundations}.

Together, these foundational concepts establish software architecture as a primary design concern and provide the basis for reasoning about architectural change, degradation, and quality-oriented design~\cite{bass2021software}.

\section{Maintainability Foundations}

Maintainability is widely recognized as a key quality attribute for long-lived software systems. It reflects the ease with which a system can be analyzed, modified, tested, and evolved. As systems increase in size and complexity, architectural decisions play a growing role in determining maintainability, often outweighing the impact of local code-level practices~\cite{bass2021software}.

Empirical studies show that maintainability evolves over time rather than remaining stable. Maintainability in open-source systems follows observable evolutionary trends, where mature systems tend to stabilize while significant degradation often coincides with major feature additions or architectural changes~\cite{molnar2020study}. These findings also indicate that targeted refactoring can mitigate maintainability degradation, highlighting the importance of deliberate architectural improvement. Industrial evidence from Microsoft further confirms that developers regularly refactor to improve maintainability, but face significant effort and risk during these activities~\cite{kim2014refactoring}.

From a measurement perspective, standardized quality models provide a common basis for evaluation. The ISO/IEC~25010 quality model defines maintainability as a composite attribute consisting of modularity, reusability, analysability, modifiability, and testability~\cite{ISO25010}. This decomposition enables maintainability to be assessed through observable structural and behavioral properties rather than subjective judgment.

At the architectural level, these sub-characteristics are strongly influenced by system decomposition, dependency structure, and clarity of component responsibilities. For example, modular designs with clear interfaces support analysability and modifiability, while low coupling and high cohesion contribute to testability and controlled change impact~\cite{ISO25010, bass2021software}. As a result, ISO/IEC~25010 provides a suitable theoretical foundation for analyzing the impact of architectural decisions on maintainability.

The Maintainability Index (MI), a composite metric combining Halstead Volume~\cite{halstead1977elements}, Cyclomatic Complexity~\cite{mccabe1976complexity}, and lines of code, has been widely adopted as a practical proxy for maintainability sub-characteristics~\cite{oman1992maintainability}. However, its validity is debated, particularly because static analysis tools that compute related quality indicators exhibit very low agreement among each other, with reported overlap of less than 0.4\% across six tools~\cite{lenarduzzi2023critical}. This motivates a multi-metric approach to maintainability assessment, where MI is complemented by structural and dependency-level indicators.

\section{Pattern--Tactic Interaction}

Architectural patterns define high-level structural organization, while architectural tactics represent localized design decisions aimed at influencing specific quality attributes~\cite{bass2021software}. The interaction between patterns and tactics can be understood by examining the structural changes required to introduce a tactic into an existing architecture~\cite{harrison2010how}. These changes range from minimal modifications within existing components to substantial restructuring that alters the original pattern.

A classification of tactic-related changes helps estimate implementation effort and architectural impact. The effectiveness and cost of a tactic depend on its compatibility with the underlying architectural pattern, meaning that tactics are not universally applicable~\cite{harrison2010how}. Quality-driven architectural design further clarifies this relationship: the applicability of tactics is constrained by the structural properties of architectural patterns and their control- and data-flow semantics~\cite{kim2009qualitydriven}. Empirical evidence shows that architectural patterns are rarely implemented without modification, and most real-world systems rely on pattern--tactic combinations rather than pure pattern implementations~\cite{kassab2018software}.

Architecture erosion poses a major threat to long-term maintainability. Common causes include inappropriate architectural changes, accumulation of technical debt, and uncontrolled growth in system complexity~\cite{li2021understanding}. Various detection practices are used in industry and research, including dependency structure analysis, architecture conformance checking, architectural smell detection, and visualization techniques~\cite{li2021understanding}. Closely related to erosion, architectural drift refers to the gradual divergence between intended and implemented architecture. This divergence can be detected by comparing architectural models extracted from source code with the original design over multiple development iterations~\cite{rosik2011assessing}. Both erosion and drift represent measurable architectural degradation that directly impacts maintainability.

Empirical mining of developer discussions reveals that architectural tactics have both positive and negative effects on quality attributes. Some tactics intended to improve one quality attribute may unintentionally harm maintainability, underscoring the need for systematic evaluation rather than relying solely on prescriptive guidance~\cite{bi2021mining}. A recent systematic mapping by M\'arquez et al.\ catalogues a large set of architectural tactics together with their primary quality attribute impacts and provides a reusable structure for tactic-based reasoning~\cite{marquez2022architectural}. Bogner et al.\ further show that modifiability tactics, in particular coupling reduction, are systematically addressed through service-oriented design patterns in both SOA and microservice architectures~\cite{bogner2019modifiability}. Rahmati and Tanhaei point out that misapplication of architectural patterns can lead to a measurable increase in maintenance effort, again emphasizing that patterns and tactics must be selected with respect to context~\cite{rahmati2021ensuring}.

\section{Maintainability Tactics}
\label{sec:maintainability_tactics}

This section summarizes the architectural tactics implemented in the LLM-based pipeline. The catalog is scoped to four tactics that are applicable to Python backend repositories and implementable through code-level transformations. The scope is narrower than full architectural tactic taxonomies such as those by Bass et al.~\cite{bass2021software} and M\'arquez et al.~\cite{marquez2022architectural}; tactics requiring runtime infrastructure changes (e.g., logging, monitoring), organizational processes (e.g., code review checklists), or deployment-level restructuring (e.g., microservice extraction) are deliberately excluded as they fall outside the code-level transformation scope of this investigation.

\subsection{Implemented Tactics}

\begin{itemize}
    \item \textbf{Decomposability (Modularization).} Split monolithic files or modules into smaller cohesive units to improve analysability and modularity. This is the most fundamental tactic and the most frequently applicable: many Python backend repositories contain large files with mixed responsibilities. Decomposability directly reduces Halstead Volume and Cyclomatic Complexity, the primary inputs to the Maintainability Index~\cite{rahmati2021ensuring}.

    \item \textbf{Localized Modification (Encapsulation).} Encapsulate change-prone logic to reduce the ripple effect of future modifications. This tactic targets files with multiple unrelated responsibilities and extracts them into focused modules with narrow interfaces. It is closely related to information hiding and is most effective in repositories with identifiable domain boundaries~\cite{kim2009qualitydriven}.

    \item \textbf{Reduced Coupling.} Reduce inter-module dependencies through abstraction, mediation, or dependency inversion. This tactic targets modifiability by introducing intermediate abstractions that decouple modules, preventing unplanned dependency propagation~\cite{kim2009qualitydriven}.

    \item \textbf{Deferred Binding Time (Late Binding).} Externalize configuration or defer concrete implementation binding to improve flexibility and substitutability. Typical applications include introducing dependency injection, configuration files, or factory methods that postpone implementation decisions~\cite{rahmati2021ensuring}.
\end{itemize}

The LLM selects among these four tactics based on the detected architectural style, baseline structural metrics, and identified deficiencies such as high fan-out or monolithic files with low MI scores. Tactic selection is inherently constrained by the pattern--tactic fit: a tactic that is cheap in one architectural pattern may be expensive or infeasible in another~\cite{harrison2010how}. Some tactics may improve one maintainability sub-characteristic while degrading another, so evaluation must consider both positive and negative impacts~\cite{bi2021mining}.

\section{Maintainability Assessment Methods}

Assessing maintainability at the architectural level requires operational metrics that connect abstract quality attributes to concrete system properties. A structured framework decomposes maintainability into measurable architectural factors aligned with sub-characteristics such as testability, changeability, analyzability, and stability~\cite{rahmati2021ensuring}. These factors translate architectural principles, such as low coupling, high cohesion, and decomposability, into actionable evaluation criteria~\cite{bass2021software}.

Static analysis tools are widely used to automate maintainability assessment. Comparative studies of popular tools reveal large differences in detection capabilities, precision, and rule coverage~\cite{lenarduzzi2023critical}. Extremely low agreement among tools indicates that assessment results depend strongly on tool selection, which highlights the need to combine multiple sources of evidence when evaluating architectural changes. As a practical consequence, recent empirical studies use the Maintainability Index together with structural metrics such as fan-in, fan-out, and dependency cycles to triangulate maintainability changes after refactoring.

\section{LLMs for Code Refactoring and Architectural Tasks}

Recent advances in large language models have enabled new forms of automated software maintenance, ranging from local code refactoring to repository-level architectural reasoning. Empirical evaluations show that LLM-based refactoring can preserve program behavior while improving maintainability-related attributes such as readability, analysability, and modifiability~\cite{depalma2024exploring}. Common refactoring actions identified in these studies include renaming, decomposition, code simplification, and structural cleanup. These recurring patterns are strongly associated with improved maintainability, suggesting that large language models can serve as effective tools for applying maintenance-related refactorings at scale~\cite{depalma2024exploring}.

\subsection{Capabilities and Limitations of LLM-Based Refactoring}

Multi-agent frameworks such as MANTRA report high success rates in producing compilable, test-passing refactored code through contextual retrieval-augmented generation and multi-agent LLM collaboration~\cite{xu2025mantra}. However, empirical analysis of AI coding agents reveals a significant scale limitation: agentic refactorings are dominated by low-level, localized edits such as renaming, type changes, and small extractions, while higher-level design changes remain comparatively rare~\cite{horikawa2025agentic}. This is particularly problematic for architectural tactics, which by nature require coordinated changes across multiple modules.

Research into prompt engineering shows that while LLMs can identify a large fraction of refactoring opportunities, their performance in complex synthesis tasks remains limited~\cite{liu2025exploring}. Piao et al.\ explored bridging human expertise with LLM-based machine understanding for refactoring and found that structured prompts that include architectural context significantly improve output quality~\cite{piao2025refactoring}. In security-related tactic implementation, Shokri and Khatchadourian reported high syntactic correctness but a sharp drop in semantic correctness, primarily due to the difficulty of coordinating inter-procedural changes across multiple classes~\cite{shokri2024ipsynth}. A recent systematic literature review of LLM-based refactoring confirms that architecture-level transformations remain an open challenge, with the majority of LLM-based studies addressing only code-level transformations~\cite{martinez2025refactoring}.

LLMs have also been applied to maintenance-oriented refactoring of Python ML projects, where multi-agent systems were used to perform a sequence of refactoring steps; the resulting work mostly focused on code-level patterns rather than full architectural reasoning~\cite{pucho2025refactoring}. A recent survey on generative AI for software architecture explicitly identifies architecture detection and tactic-based refactoring as key open areas for future research~\cite{esposito2025generative}.